\section{Proof of the contraction criterion}
\label{sec:appendix-contraction}

\begin{proof}[Proof of \Cref{thm:contraction}]
Let
\[
  L_n(u)
  =\frac{\mathrm dP_u^{(n)}}{\mathrm dP_{\udag}^{(n)}}
      (Y^{(n)}),
  \qquad
  D_n=\int L_n(u)\,\muz(\mathrm du)
\]
be the likelihood ratio and the posterior denominator.  The usual
Bayesian denominator lemma, applied to
\eqref{eq:kl-neighborhood}--\eqref{eq:contraction-mass}, gives
\begin{equation}
  P_{\udag}^{(n)}
  \left(
    D_n<e^{-(c+2)n\epsilon_n^2}
  \right)\longrightarrow0 .
  \label{eq:denominator-event}
\end{equation}
The lemma follows by integrating the log-likelihood ratios over
$\mathcal K_n$ and applying Chebyshev's inequality to the variance
bound; it is the standard denominator step in posterior-contraction
proofs \citep[Chapter~8]{ghosal2017}.

Set
\[
  A_n
  =\{u\in\mathcal F_n:d_n(u,\udag)>M\epsilon_n\}.
\]
Fubini's theorem and change of measure give
\begin{align*}
  \mathbb E_{\udag}^{(n)}
  \left[
    (1-\varphi_n)\int_{A_n}L_n(u)\,\muz(\mathrm du)
  \right]
  &=
  \int_{A_n}P_u^{(n)}(1-\varphi_n)\,\muz(\mathrm du)\\
  &\leq e^{-(c+4)n\epsilon_n^2}.
\end{align*}
On the event in which the lower bound complementary to
\eqref{eq:denominator-event} holds, Markov's inequality therefore
shows that the posterior mass of $A_n$, apart from the type-I test
term, converges to zero.  Indeed, the ratio of the displayed
expectation to the denominator lower bound is at most
$e^{-2n\epsilon_n^2}$.

For the complement of the sieve, Fubini gives
\[
  \mathbb E_{\udag}^{(n)}
  \int_{\mathcal F_n^c}L_n(u)\,\muz(\mathrm du)
  =\muz(\mathcal F_n^c)
  \leq e^{-(c+4)n\epsilon_n^2},
\]
and the same denominator argument applies.  Finally,
$P_{\udag}^{(n)}\varphi_n\to0$ by
\eqref{eq:contraction-type1}.  Combining the test term, the two
numerator bounds, and \eqref{eq:denominator-event} proves
\eqref{eq:contraction-conclusion}.
\end{proof}
